% Phase 7 research supplement. Author: Wenyu Huang, UCSD.
% This is a complete proof candidate, not a claim of a new literature theorem.
% The original project and Phase 6 have not been modified.
\section{A smooth bounded-slope generator with unbounded excess hierarchy price}
\label{sec:smooth-counterexample}

\paragraph{Scope.}
The divergent-endpoint-slope theorem in \texttt{manuscript/general\_losses.tex}
is a sufficient condition, not a characterization.  We prove that its
endpoint-slope hypothesis is not necessary.  The conclusion below concerns
\emph{excess} Bayes risk.  It supplies neither the optimal universal
four-state logarithmic-loss coefficient nor a new total-risk obstruction.

Write $\psi(t)=t\log t+(1-t)\log(1-t)$.  For a finite source with positive
masses $p_x$ and binary posteriors $\eta_x$, put
\[
 D_\phi(\mathcal P)=\sum_xp_x\phi(\eta_x)
 -\sum_{C\in\mathcal P}p_C\phi\left(\sum_{x\in C}p_x\eta_x/p_C\right).
\]
The flat benchmark $D_{\phi,k}^*$ minimizes this quantity over all
$k$-cell partitions.  A deterministic four-state hierarchy has nested
partitions at budgets two and three, and its price is the minimum over
hierarchies of $\max_{k=2,3}D_\phi(\mathcal H_k)/D_{\phi,k}^*$.
For an encoder $X\to Z$ independent of $Y$ conditional on $X$, define
\[
 D_\phi(Z)=\mathbb E\phi(\eta_X)
             -\mathbb E\phi(\mathbb E[\eta_X\mid Z]).
\]
For stochastic hierarchies use $Y-X-Z_3-Z_2$ with output budgets three
and two and the same flat denominators.  Zero-probability encoder outputs
are omitted from conditional expectations.

\begin{theorem}[Finite endpoint slopes do not characterize bounded price]
\label{thm:smooth-bounded-slope}
There exists one fixed function $\phi\in C^\infty([0,1])$ such that:
\begin{enumerate}
\item $\phi(t)=\phi(1-t)$, $\phi''(t)>0$ for every $0<t<1$, and
      $0\leq\phi''(t)\leq1$ on $[0,1]$;
\item $\phi'(0)$ and $\phi'(1)$ are finite, while
      $\phi^{(r)}(0)=\phi^{(r)}(1)=0$ for every integer $r\geq2$;
\item there is a sequence of full-support four-by-two laws, each with
      four distinct posteriors, whose deterministic excess hierarchy
      prices obey
      \[
       \rho_{\phi,\mathrm{det},j}\sim
       \frac{\sqrt j}{2\sqrt{\log2}}\longrightarrow\infty.
      \]
      Their stochastic excess hierarchy prices likewise obey
      \[
       \rho_{\phi,\mathrm{stoch},j}\sim
       \frac{\sqrt j}{4\sqrt{\log2}}\longrightarrow\infty.
      \]
\end{enumerate}
The generator is strictly convex, but not uniformly strongly convex on
the closed interval.  Its associated binary strictly proper losses can
be chosen smooth and valued in $[0,1]$.
\end{theorem}

\begin{proof}
We first construct a smooth positive curvature containing disjoint exact
copies of truncated Shannon curvature.  Fix a smooth function
$\vartheta:\mathbb R\to[0,1]$ equal to zero on $(-\infty,0]$ and one on
$[1,\infty)$.  For integers $j\geq2$, set
\[
 a_j=2^{-j-2},\quad l_j=2^{-j-3},\quad
 I_j=(a_j,a_j+l_j),\quad e_j=e^{-2j},\quad \lambda_j=e^{-4^j}.
\]
The intervals $I_j$ are disjoint, lie below $1/4$, and accumulate only
at zero.  Define
\[
 \chi_j(t)=\vartheta\left(\frac{t-e_j/2}{e_j/2}\right)
           \vartheta\left(\frac{1-e_j/2-t}{e_j/2}\right).
\]
Thus $\chi_j=1$ on $[e_j,1-e_j]$, it vanishes outside
$(e_j/2,1-e_j/2)$, and, for every fixed integer $m\geq0$,
$\|\chi_j^{(m)}\|_\infty\leq C_m e_j^{-m}$.

Use the positive symmetric baseline
\[
 b(x)=\exp\{-1/[x(1-x)]\}\quad(0<x<1),\qquad b(0)=b(1)=0.
\]
For $x\in I_j$, write $t=(x-a_j)/l_j$ and set
\[
 g(x)=(1-\chi_j(t))b(x)
       +\chi_j(t)\frac{\lambda_j}{l_j^2t(1-t)}.
 \tag{1}\label{eq:curvature-patch}
\]
On $1-I_j$ define $g(x)=g(1-x)$; elsewhere set $g=b$.
The cutoff has support strictly inside each interval, so this is smooth
at the endpoints of every individual interval and positive throughout
$(0,1)$.  The only smoothness issue is the accumulation at zero and one.

Here are explicit estimates for that issue.  On the support of $\chi_j$,
\[
 \left|\frac{d^r}{dt^r}\frac1{t(1-t)}\right|
 \leq C_r e_j^{-r-1}.
\]
Leibniz' rule and rescaling therefore bound the $m$th derivative of the
second summand in (1) by
\[
 C_m\lambda_j l_j^{-m-2}e_j^{-m-1}
 \leq C_m\exp\{-4^j+C_mj\}.
 \tag{2}\label{eq:target-decay}
\]
On $I_j$, derivatives of $b$ are bounded by a polynomial in $2^j$
times $\exp\{-c2^j\}$ for an absolute $c>0$.  Applying Leibniz'
rule once more shows that the $m$th derivative of the first summand in
(1) is at most
\[
 C_m\exp\{-c2^j+C_mj\}.
 \tag{3}\label{eq:baseline-decay}
\]
For clarity, differentiating $b$ $r$ times yields $b$ times a rational
function with poles only at zero and one; on $I_j$ its magnitude is at
most $C_r2^{C_rj}$, while $x\leq3\cdot2^{-j-3}$ gives the exponential
factor in (3).  Each cutoff derivative contributes at most
$C_r l_j^{-r}e_j^{-r}=\exp\{O_r(j)\}$.

The bounds (2)--(3) vanish faster than $a_j^N$ for every fixed $m,N$.
Outside the intervals the baseline has the same flatness property.
Consequently all derivatives of $g$ on $(0,1)$ tend to zero faster than
any power of $x$ at zero.  Extending $g^{(m)}(0)=0$ for all $m$ is
consistent with differentiation: the difference quotient of the
extended $g^{(m)}$ tends to zero because $g^{(m)}(x)=o(x)$.
Induction proves $g\in C^\infty([0,1])$, flat at zero.  Reflection proves
the same at one.

Let $M=\max_{[0,1]}g>0$ and define $\phi$ by
\[
 \phi''=g/M,\qquad \phi(1/2)=\phi'(1/2)=0.
 \tag{4}\label{eq:integrate-curvature}
\]
It is smooth, symmetric, strictly convex, and has $0\leq\phi''\leq1$.
Strict convexity follows because every nonempty interior interval has
strictly positive curvature integral.  Its endpoint derivatives are
finite by (4); its derivatives of order at least two vanish there.

On the closed interval
$J_j=[a_j+l_je_j,a_j+l_j(1-e_j)]$, (1) and (4) imply
\[
 \phi(x)=\frac{\lambda_j}{M}
          \psi\left(\frac{x-a_j}{l_j}\right)+A_jx+B_j
 \tag{5}\label{eq:exact-local-copy}
\]
for some constants $A_j,B_j$.  This follows by matching second
derivatives and integrating twice; the constants need not be evaluated.

Now take the already established logarithmic-loss four-state family
at parameter $L=j$:
\[
 w_j=e^{-j},\quad c_j=(1-2w_j)/2,\quad
 \delta_j=\sqrt{\log2/j},
\]
with masses $(c_j,w_j,w_j,c_j)$ and original posteriors
\[
 (q_A,q_C,q_D,q_B)=(e_j,\delta_j,1-\delta_j,1-e_j).
\]
For all sufficiently large $j$, these are four distinct points in
$[e_j,1-e_j]$.  Keep the masses and replace each posterior by
\[
 \eta_x=a_j+l_jq_x.
 \tag{6}\label{eq:transformed-posteriors}
\]
Every new posterior is strictly between zero and one, so the resulting
four-by-two joint law has full support.

Every cell posterior, and every conditional posterior of an arbitrary
stochastic encoder of $X$, is a convex combination of the four values
in (6).  It therefore belongs to $J_j$.  Equation (5), affinity of the
posterior transformation, and cancellation of affine terms give the
\emph{exact} identities
\[
 D_\phi(\mathcal P)=\frac{\lambda_j}{M}D_\psi(\mathcal P),\qquad
 D_\phi(Z)=\frac{\lambda_j}{M}D_\psi(Z)
 \tag{7}\label{eq:all-encoder-transfer}
\]
for every partition and every stochastic encoder.  Hence the flat
optima and all attainable two-level distortion vectors are scaled by
the same positive number.  In particular all hierarchy ratios, including
their minimax values, are unchanged.  The deterministic excess
asymptotic in the four-state theorem immediately proves the stated
deterministic conclusion.

For completeness, the existing stochastic theorem is stated for total
logarithmic risk, but it implies the same excess asymptotic for this
family.  Write $R_k^*=B_{0,j}+D_{\psi,k}^*$ and
$b_j=\max_{k=2,3}B_{0,j}/D_{\psi,k}^*$.  The established cost estimates
give $B_{0,j}=o(w_j)$, $D_{\psi,3}^*\sim2w_j\log2$ and
$D_{\psi,2}^*\sim2w_j\delta_jj$, so $b_j\to0$.  For every admissible
chain its excess and total ratios satisfy
\[
 \frac1{1+b_j}\max_k\frac{D_\psi(Z_k)}{D_{\psi,k}^*}
 \leq\max_k\frac{B_{0,j}+D_\psi(Z_k)}{B_{0,j}+D_{\psi,k}^*}
 \leq\max_k\frac{D_\psi(Z_k)}{D_{\psi,k}^*}.
\]
The right inequality uses $D_\psi(Z_k)\geq D_{\psi,k}^*$; flat
stochastic encoders do not improve a concave Bayes-risk objective,
by the existing random-function argument.  Infimizing proves equality
of the stochastic total and excess prices up to $1+o(1)$, and (7)
transfers the result to $\phi$.

Finally put $h(q)=\phi(0)-\phi(q)$ and define the proper losses
\[
 \ell(q,1)=h(q)+(1-q)h'(q),\qquad
 \ell(q,0)=h(q)-qh'(q).
\]
Their expected regret at truth $p$ is
$\phi(p)-\phi(q)-\phi'(q)(p-q)$, strictly positive for $p\ne q$.
Convexity at the endpoints makes both losses nonnegative.  Moreover
$\ell(q,1)'=-(1-q)\phi''(q)\leq0$ and
$\ell(q,0)'=q\phi''(q)\geq0$.  Their largest endpoint values equal
$-\phi'(0)=\phi'(1)=\int_0^{1/2}\phi''(t)dt\leq1/2$.
Thus both smooth losses lie in $[0,1]$ as claimed.
\end{proof}

\paragraph{Why this does not give a total-risk counterexample.}
For the displayed proper losses, $h(q)=cq+o(q)$ at zero, with
$c=-\phi'(0)>0$.  All posteriors in (6) lie in $[a_j,3a_j/2]$.
The full-information Bayes risk $B_{\phi,j}=\sum_xp_xh(\eta_x)$ is
therefore at least $ca_j/2$ eventually.  On the other hand (7) and
the entropy bound imply, for every encoder,
$0\leq D_\phi(Z)\leq(\lambda_j/M)\log2$.  Hence every ratio of total
risk to a flat total-risk optimum is bounded above by
\[
 1+\frac{2\lambda_j\log2}{Mca_j}=1+o(1).
\]
In this construction all total-risk hierarchy prices tend to one.  The
unbounded excess ratio occurs at vanishing absolute distortion and is
not a claim of large absolute loss or practical learning difficulty.

\paragraph{Novelty boundary.}
The proof imports the existing four-state lower bound and adds a smooth
local-embedding argument.  Bregman affine invariance, proper-score
representations, smooth cutoff functions, and bounded-curvature
comparison are established tools.  The precise fixed-generator
counterexample is a candidate boundary result; no claim of priority is
made without the accompanying primary-source audit.
